\section{Related Work}
\label{sec:rw}

\subsection{Classical and neural retrieval}
\label{sec:rw-classical}

The two lexical baselines evaluated in this paper follow the standard formulations: TF-IDF term-weighting as systematised by Salton \& Buckley (\citep{salton1988termweighting}) and the BM25 probabilistic relevance framework (\citep{robertson2009probabilistic}). Their neural successors replace sparse term matching with dense representations: DPR-style bi-encoders encode query and document independently into a shared vector space (\citep{karpukhin2020dense}), while late-interaction models such as ColBERT retain per-token contextualised embeddings and score with token-level interactions, recovering much of the lexical precision that pure bi-encoders lose (\citep{khattab2020colbert}). Our Dense-IR and ColBERT-like baselines are canonical representatives of these two families, and --- consistent with the literature --- the late-interaction model is the stronger of the two on every metric in our leaderboard (\Cref{sec:results-sota}). None of these systems maintains state across turns, which motivates the memory mechanisms below.

\subsection{Agentic and conversational IR}
\label{sec:rw-agentic}

The evaluation setting in this paper --- retrieval over multi-turn conversational sessions --- sits at the intersection of two active lines of work. On the evaluation side, TREC's conversational-assistance track established multi-turn, context-dependent retrieval as a standard task (\citep{dalton2020cast}) , and FIRE has run sustained shared tasks with a strong Indian-language and code-mixed emphasis whose conventions our benchmarks mirror (FIRE proceedings \citep{fire2025proceedings}) . On the systems side, ``agentic IR'' has recently been articulated as a shift from static query--document matching to goal-directed, iterative retrieval in which an agent reasons, retrieves, and adapts over multiple steps (\citep{zhang2024agentic}) , with follow-up work on reasoning-augmented deep research agents (\citep{zhang2025websearch})  and on simulating search users for realistic multi-turn evaluation (\citep{zhang2024usimagent}) . MEIRA belongs to this paradigm: each conversation is a session over which the system accumulates memory and refines its retrievals; we differ from the position papers in the concrete mechanisms --- an episodic memory bank, a decay schedule, and an attribution head --- and in measuring them (XAIR@K, MDS).

\subsection{Memory-augmented LLMs and retrieval}
\label{sec:rw-memory}

Our episodic memory bank draws on two influential lines of work on giving LLM-based systems long-term memory. MemGPT (\citep{packer2023memgpt})  treats the model as an operating system with hierarchical memory tiers, paging information between a main context and external archival storage via explicit tool calls; MEIRA's 64-slot bank is a retrieval-specialised version of this idea, persisting \emph{retrieval states} rather than raw dialogue. Generative Agents (\citep{park2023generative})  introduced a memory stream with recency, importance and relevance scoring that determines what an agent recalls --- a direct antecedent of our temporally-decayed memory. Closest to our decay mechanism is MemoryBank (\citep{zhong2024memorybank}) , which grounds memory consolidation and forgetting in the Ebbinghaus forgetting curve so that older, less-salient memories decay over time. MEIRA operationalises the same intuition for retrieval: temporal decay down-weights stale memory slots, and the ablation (\Cref{sec:results-ablation}) shows it is the second-largest contributor to performance ($\Delta$F1 +0.078/+0.086) --- the empirical counterpart to the theoretical case for forgetting made in these works.

\subsection{Explainable and trustworthy IR}
\label{sec:rw-xai}

The XAI attribution head and the XAIR@K metric position MEIRA within the explainable IR (ExIR) literature. \citet{anand2022explainable}  survey the field, distinguishing transparent-by-design from post-hoc methods and noting the absence of standardised evaluation of explanation quality --- the exact gap XAIR@K targets by folding attribution confidence into the ranking score. Recent surveys extend this picture to LLM-based systems, cataloguing explanation generation and human-centred trust evaluation for generative models (\citep{zhao2024explainability})  and LLM-driven explainable AI more broadly (\citep{bilal2025llms}) . Relative to this literature our contribution is measurement-oriented: rather than proposing a new attribution method, we propose a \emph{metric} that makes explainability part of the leaderboard, and we show (\Cref{sec:results-ablation}) that the attribution head is the sole driver of XAIR@10 ($\Delta$XAIR@10 +0.894/+0.886) --- i.e. that the metric responds exactly to the mechanism it is designed to measure.

\subsection{Hard negatives, evaluation practice, and the proposed metrics}
\label{sec:rw-hardneg}

Two further strands inform the evaluation design. First, hard-negative construction: sibling-topic negatives follow the established finding that hard negatives --- topically confusable distractors --- make ranking tasks non-trivial and drive separation between lexical and neural systems (Karpukhin et al., 2020 ; \citet{wischounig2026negative} ); our corpus statistics confirm this (67.1\% / 63.9\% of negatives are hard). Second, evaluation methodology: LLM-as-judge approaches meta-evaluate judges' agreement with humans (\citep{li2024llmsjudges}) , and answerability-aware metrics argue that graded relevance alone under-describes retrieval quality (\citep{farzi2024exam}) . Our two proposed metrics sit naturally in this conversation: XAIR@K argues that a \emph{correct-but-unexplainable} hit is worth less than a \emph{correct-and-explained} one, and MDS argues that a system's \emph{utilisation of its own memory} is a measurable, reportable property --- both are checks on behaviours that standard graded-relevance metrics are blind to.

\subsection{Positioning}
\label{sec:rw-positioning}

In summary: MEIRA is an agentic retrieval system (\Cref{sec:rw-agentic}) whose defining mechanisms are memory with forgetting (\Cref{sec:rw-memory}) and explainability (\Cref{sec:rw-xai}), evaluated on benchmarks that stress hard negatives and cross-lingual multi-turn retrieval (\Cref{sec:rw-classical,sec:rw-agentic,sec:rw-hardneg}). Its novelty is not any single mechanism --- all exist in the literature --- but their \emph{coupling inside one retrieval agent} and, we argue, the \emph{measurement} of them: XAIR@K and MDS make explainability and memory utilisation first-class citizens of the leaderboard. The evaluation (\Crefrange{sec:results}{sec:rob}) then holds MEIRA to the same statistical rigour as any leaderboard claim: paired tests, multiplicity correction, and a threshold sweep.
